\documentclass[11pt,a4paper]{article}

\usepackage{main}
\usepackage{cours}

\RequirePackage{tikz}
\usetikzlibrary{tikzmark, arrows.meta}

\renewcommand{\typedocument}{Fiche}
\renewcommand{\classe}{}
\renewcommand{\titre}{Fiche - Grandeurs et unités}

\begin{document}

\creertitre{\titre}

Une \important{grandeur physique} est une propriété d'un objet ou d'un phénomène qui peut être mesurée ou calculée.

On l'exprime avec une \textbf{valeur numérique} et une \textbf{unité de mesure}.

\exemple{Une table a une
\tikzmarknode{grandeur}{masse}
\tikzmarknode{symbole}{$m$}
$=$
\tikzmarknode{valeur}{$10$}
\tikzmarknode{unite}{\si{\kilo\gram}}%
.}

\begin{tikzpicture}[remember picture, overlay,
    >=Stealth, every node/.style={font=\small}]
  \draw[->] ([yshift=-4pt]grandeur.south)
      -- ++(-1.1,-0.8) node[below] {grandeur};
  \draw[->] ([yshift=-4pt]symbole.south)
      -- ++(-0.3,-0.8) node[below] {symbole};
  \draw[->] ([yshift=-4pt]valeur.south)
      -- ++(0.3,-0.8) node[below] {valeur};
  \draw[->] ([yshift=-2pt]unite.south)
      -- ++(1.1,-0.8) node[below] {unité};
\end{tikzpicture}

\vspace{1em}
\partie{Grandeurs physiques}
\vspace*{-0.7em}

\renewcommand{\arraystretch}{1.6}
\begin{center}
\small
\begin{tabular}{|l|c|l|l|l|}
\hline
\rowcolor{headergrey}
\textbf{Grandeur} & \textbf{Symbole} & \textbf{Unité SI (symbole)} & \textbf{Autres unités} & \textbf{Appareil de mesure} \\
\hline
Distance & $\ell$ & mètre (\si{\metre}) & kilomètre (\si{\km}) & Règle, mètre \\
\hline
Durée & $t$ & seconde (\si{\second}) & minute (\si{\minute}), heure (\si{\hour}) & Chronomètre \\
\hline
Vitesse & $v$ & mètre par seconde (\si{\metre/\second}) & kilomètre par heure (\si{\km/\hour}) & \\
\hline
Fréquence & $f$ & hertz (\si{\hertz}) & & \\
\hline
Masse & $m$ & kilogramme (\si{\kg}) & gramme (\si{\g}), tonne (\si{\tonne}) & Balance \\
\hline
Volume & $V$ & mètre cube (\si{\cubic\metre}) & litre (\si{\litre}), millilitre (\si{\milli\litre}) & Éprouvette graduée \\
\hline
Masse volumique & $\rho$ & kilogramme par mètre cube (\si{\kg/\cubic\metre}) & \si{\g/\milli\litre}, \si{\g/\centi\cubic\metre} & \\
\hline
pH & pH & Pas d'unité & & Papier pH \\
\hline
Poids & $P$ & newton (\si{\newton}) & & Dynamomètre \\
\hline
Température & $T$ & kelvin (\si{\kelvin}) & degrés Celsius (\si{\degreeCelsius}) & Thermomètre \\
\hline
Tension électrique & $U$ & volt (\si{\volt}) & & Voltmètre \\
\hline
Intensité électrique & $I$ & ampère (\si{\ampere}) & & Ampèremètre \\
\hline
Résistance électrique & $R$ & ohm (\si{\ohm}) & & \\
\hline
Puissance & $P$ & watt (\si{\watt}) & & \\
\hline
Énergie & $E$ & joule (\si{\joule}) & kilowatt-heure (\si{\kilo\watt\hour}) & \\
\hline
\end{tabular}
\end{center}

\begin{minipage}[t]{0.55\linewidth}
    \partie{Préfixes des unités de mesure}
    \vspace*{-0.7em}

    \renewcommand{\arraystretch}{1.6}
    \begin{center}
    \small
    \begin{tabular}{|p{1.75cm}|p{1.75cm}|p{1.75cm}|p{3.5cm}|}
    \hline
    \rowcolor{headergrey}
    \textbf{Préfixe} & \textbf{Symbole} & \textbf{Facteur} & \textbf{Exemple} \\
    \hline
    méga & M & $\times 10^{6}$ & \SI{1}{\mega\hertz} = \SI{1000000}{\hertz} \\
    \hline
    kilo & k & $\times 10^{3}$ & \SI{1}{\km} = \SI{1000}{\metre} \\
    \hline
    hecto & h & $\times 10^{2}$ & \SI{1}{\hecto\gram} = \SI{100}{\gram} \\
    \hline
    deca & da & $\times 10^{1}$ & \SI{1}{\deca\metre} = \SI{10}{\metre} \\
    \hline
    déci & d & $\times 10^{-1}$ & \SI{1}{\deci\litre} = \SI{0.1}{\litre} \\
    \hline
    centi & c & $\times 10^{-2}$ & \SI{1}{\cm} = \SI{0.01}{\metre} \\
    \hline
    milli & m & $\times 10^{-3}$ & \SI{1}{\milli\second} = \SI{0.001}{\second} \\
    \hline
    micro & µ & $\times 10^{-6}$ & \SI{1}{\micro\gram} = \SI{0.000001}{\gram} \\
    \hline
    \end{tabular}
    \end{center}
    \renewcommand{\arraystretch}{1}
\end{minipage}
\hfill
\begin{minipage}[t]{0.4\linewidth}
    \partie{Conversions courantes}

    \begin{itemize}
        \item $\SI{1}{\tonne} = \SI{1000}{\kilo\gram}$
        \item $\SI{1}{\litre} = \SI{0.001}{\cubic\metre} \leftrightarrow \SI{1}{\cubic\metre} = \SI{1000}{\litre}$
        \item $\SI{1}{\deci\cubic\meter} = \SI{1}{\litre}$
        \item $\SI{1}{\centi\cubic\meter} = \SI{1}{\milli\litre}$
        \item $\SI{1}{\g/\litre} = \SI{1}{\kg/\cubic\metre}$
        \item $\SI{1}{\hour} = \SI{3600}{\second} \leftrightarrow \SI{1}{\second} = \frac{1}{3600} \SI{}{\hour}$
        \item $\SI{1}{\meter/\second} = \SI{3.6}{\kilo\meter/\hour} \leftrightarrow \SI{1}{\kilo\meter/\hour} = \frac{1}{3.6} \SI{}{\meter/\second}$
    \end{itemize}
\end{minipage}

\end{document}